% Proof-ready insertion after eq:prime-tail-target in local-energy.tex.
% Uses the classical Davenport estimate already invoked in sec:sieve.
\begin{proposition}[Prime-tail cancellation to logarithmic order]
\label{prop:prime-tail-logarithmic}
For \(y=(\log X)^2\), and every fixed \(A>0\),
\[
 T_y(X)=-c_yX+O_A\!\left(\frac{X}{(\log X)^A}\right).
\]
In particular, \(T_y(X)/(c_yX)\to-1\).
\end{proposition}
\begin{proof}
The classical uniform additive estimate of Davenport, in the
Liouville form recalled in Section~\ref{sec:sieve} and quoted in
\cite{mrt}, gives at additive frequency zero
\[
 L(t)\ll_A \frac{t}{(\log(2t))^A}
 \qquad(t\ge1).
\]
Partial summation gives the absolutely convergent identity
\[
 B(X)=\frac1X\int_1^\infty L(t)e^{-t/X}\dd t.
\]
For \(X\ge4\), the contribution from \(1\le t\le\sqrt X\)
has modulus at most \(X^{-1}\int_1^{\sqrt X}t\dd t\le1/2\).
For \(t\ge\sqrt X\), use
\(\log(2t)\ge\tfrac12\log X\) to bound the remaining part by
\[
 \frac{C_A2^A}{X(\log X)^A}
     \int_0^\infty te^{-t/X}\dd t
 =\frac{C_A2^AX}{(\log X)^A}.
\]
Consequently \(B(X)\ll_A X/(\log X)^A\).
By Theorem~\ref{thm:finite-prime} and
\eqref{eq:prime-tail-exact},
\[
 T_y(X)+c_yX
   =B(X)-\bigl(B_y(X)-c_yX\bigr)
   =O_A\!\left(\frac{X}{(\log X)^A}\right)
        +O(X^{1/2+o(1)}).
\]
For each fixed \(A\), the last error is absorbed in the first.
Finally \(c_y\ge(\log X)^{-4}\) gives
\(T_y(X)/(c_yX)=-1+O_A((\log X)^{4-A})\).
Take any \(A>4\).
\end{proof}

This proves the main-term cancellation in the actual exponential
kernel. It is a consequence of classical additive cancellation,
not a new zero-free theorem. Its error is larger than the
square-root error requested in \eqref{eq:prime-tail-target}.

\begin{proposition}[The precise power-saving endpoint]
\label{prop:prime-tail-rh-equivalence}
For \(y=(\log X)^2\), the estimate
\[
 T_y(X)+c_yX\ll_\varepsilon X^{1/2+\varepsilon}
 \quad\text{for every }\varepsilon>0
\]
is equivalent to RH.
\end{proposition}
\begin{proof}
The exact identity and Theorem~\ref{thm:finite-prime} give
\[
 T_y(X)+c_yX-B(X)=O(X^{1/2+o(1)}).
\]
Hence the asserted family of bounds holds if and only if
\(B(X)\ll_\varepsilon X^{1/2+\varepsilon}\) for every
\(\varepsilon>0\). The equivalence of the latter condition
with RH is Proposition~\ref{prop:smoothed}.
\end{proof}

% Suitable following the existing Abel double-zero discussion.
For clarity, the double zero in \eqref{eq:prime-tail-double-zero}
does not exclude nonreal poles. For fixed \(y\), each finite
Euler factor of \(c_y(s)\) is holomorphic and nonzero on
\(\operatorname{Re}s>0\). If \(\rho\) were a zeta zero of
order \(r\) with \(\operatorname{Re}\rho>1/2\), then
\eqref{eq:prime-tail-abel} would have the leading Laurent term
\[
 H_y(s)=
 \frac{(r!)^2\zeta(2\rho)}
      {c_y(\rho)(\zeta^{(r)}(\rho))^2}
       (s-\rho)^{-2r}
       +O((s-\rho)^{1-2r}).
\]
The coefficient is nonzero. The local cancellation at one
therefore supplies no contradiction to such a zero; excluding
it requires the global power-saving estimate.
